\begin{helper}
Let \(V\) be a finite set and let \(\nu\) be a probability law on pairs
\((A,\pi)\), where \(A\subseteq V\) and \(\pi:V\to\mathbb R\).  Fix
\(p\in[0,1]\).  Suppose that the activation atom of every set \(A_0\subseteq V\)
has mass
\[
\nu\{A=A_0\}=
\operatorname{ofReal}\!\left(p^{|A_0|}(1-p)^{|V|-|A_0|}\right),
\]
and suppose that, in the product-law sense expressed by lower rectangles, for
every threshold vector \(s\in[0,1]^V\),
\[
\nu\{A=A_0,\ 0\le \pi(v)\le s_v\hbox{ for all }v\in V\}
=
\nu\{A=A_0\}\operatorname{ofReal}\!\left(\prod_{v\in V}s_v\right).
\]
Let \(L,U\subseteq V\) be disjoint finite coordinate sets and let
\(t_v\in[0,1]\) for every \(v\in L\cup U\).  Then, for every activated set
\(A_0\subseteq V\),
\[
\nu\{A=A_0,\ 0\le \pi(v)\le1\hbox{ for all }v\in V,\ 
      \pi(v)\le t_v\hbox{ for }v\in L,\ 
      t_v<\pi(v)\hbox{ for }v\in U\}
\]
equals
\[
\operatorname{ofReal}\!\left(
p^{|A_0|}(1-p)^{|V|-|A_0|}
\prod_{v\in L}t_v
\prod_{v\in U}(1-t_v)\right).
\]
\end{helper}
\begin{proof}
Write
\[
\alpha=p^{|A_0|}(1-p)^{|V|-|A_0|}.
\]
The activation-atom hypothesis says \(\nu\{A=A_0\}=\operatorname{ofReal}(\alpha)\).
The assumptions \(0\le p\le1\) imply \(\alpha\ge0\).

For every subset \(B\subseteq U\), define a threshold vector \(s^B\) by
\[
s^B_v=
\begin{cases}
t_v,& v\in L\cup B,\\
1,& v\notin L\cup B.
\end{cases}
\]
Because \(L\cap U=\emptyset\), a coordinate of \(s^B\) is either one of the
numbers \(t_v\) with \(v\in L\cup U\), or it is \(1\).  Hence
\(0\le s^B_v\le1\) for every \(v\).  Let
\[
E_B=\{A=A_0,\ 0\le\pi(v)\le1\hbox{ for all }v\in V,\ 
        \pi(v)\le t_v\hbox{ for }v\in L\cup B\}.
\]
The event \(E_B\) is exactly the lower rectangle in the hypothesis with
threshold vector \(s^B\): on \(L\cup B\) the upper bound is \(t_v\), and off
\(L\cup B\) the upper bound is \(1\).  Therefore
\[
\nu(E_B)
=\nu\{A=A_0\}\operatorname{ofReal}\!\left(\prod_{v\in V}s^B_v\right)
\]
\[
=
\operatorname{ofReal}\!\left(
\alpha
\prod_{v\in L}t_v
\prod_{v\in B}t_v\right).
\tag{1}
\]
In the last equality we used disjointness of \(L\) and \(U\) to split
\(\prod_v s^B_v\) into the \(L\)-part, the \(B\)-part, and factors equal to
\(1\), and we used \(\operatorname{ofReal}(xy)=\operatorname{ofReal}(x)
\operatorname{ofReal}(y)\) for nonnegative real factors.

Now work inside the finite-measure cylinder \(E_\emptyset\).  It has finite
mass by (1), because \(\operatorname{ofReal}\) of a real number is never
\(\infty\).  For each \(S\subseteq U\), let \(D_S\) be the event
\[
\begin{aligned}
D_S=\{&A=A_0,\ 0\le\pi(v)\le1\hbox{ for all }v\in V,\ 
        \pi(v)\le t_v\hbox{ for }v\in L,\\
     &t_v<\pi(v)\hbox{ for }v\in S,\ 
        \pi(v)\le t_v\hbox{ for }v\in U\setminus S\}.
\end{aligned}
\]
The events \(D_S\), \(S\subseteq U\), are pairwise disjoint: if
\(u\in S\triangle T\), then one event imposes \(t_u<\pi(u)\) and the other
imposes \(\pi(u)\le t_u\).  Their union is \(E_\emptyset\), since for each
\(u\in U\) exactly one of the two alternatives
\(\pi(u)\le t_u\) and \(t_u<\pi(u)\) holds.  More generally, for each
\(B\subseteq U\),
\[
E_B=\bigsqcup_{S\subseteq U\setminus B}D_S,
\tag{2}
\]
because the extra lower-cut requirements in \(E_B\) force every coordinate of
\(B\) to lie on the \(\pi(u)\le t_u\) side, while the remaining coordinates
of \(U\setminus B\) are free to choose either side.  Since all these sets are
subsets of the finite-measure set \(E_\emptyset\), their measures are finite.
Thus we may write \(d_S\in\mathbb R_{\ge0}\) for the real number satisfying
\(\nu(D_S)=\operatorname{ofReal}(d_S)\), and finite additivity applied to
(2) gives
\[
\nu(E_B)=\operatorname{ofReal}\!\left(\sum_{S\subseteq U\setminus B}d_S\right).
\tag{3}
\]
Set
\[
r_B=\alpha\prod_{v\in L}t_v\prod_{v\in B}t_v .
\]
Equation (1) says \(\nu(E_B)=\operatorname{ofReal}(r_B)\), and \(r_B\ge0\).
Comparing this with (3) and using injectivity of
\(\operatorname{ofReal}\) on nonnegative real numbers gives
\[
\sum_{S\subseteq U\setminus B}d_S=r_B.
\tag{4}
\]

Let \(D_U\) be the desired event, namely the event with all coordinates in
\(U\) on the strict side \(t_v<\pi(v)\).  We determine \(d_U\) from the
numbers \(r_B\) by finite Möbius inversion.  Starting from (4),
\[
\sum_{B\subseteq U}(-1)^{|B|}
  \sum_{S\subseteq U\setminus B}d_S
=
\sum_{S\subseteq U}d_S
  \sum_{B\subseteq U\setminus S}(-1)^{|B|}.
\]
The inner alternating sum is \((1-1)^{|U\setminus S|}\), so it is \(0\) unless
\(S=U\), and it is \(1\) when \(S=U\).  Hence
\[
d_U=
\sum_{B\subseteq U}(-1)^{|B|}r_B.
\tag{5}
\]
Substituting the definition of \(r_B\), this becomes
\[
d_U=
\alpha\prod_{v\in L}t_v
\sum_{B\subseteq U}(-1)^{|B|}\prod_{v\in B}t_v.
\]
The remaining finite algebra identity is the expansion
\[
\prod_{v\in U}(1-t_v)
=\sum_{B\subseteq U}(-1)^{|B|}\prod_{v\in B}t_v,
\]
obtained by choosing either the \(1\)-term or the \(-t_v\)-term in each
factor.  Therefore
\[
d_U=\alpha\prod_{v\in L}t_v\prod_{v\in U}(1-t_v).
\]

Finally, \(\nu(D_U)=\operatorname{ofReal}(d_U)\).  Replacing \(\alpha\) by
the activation atom
\(p^{|A_0|}(1-p)^{|V|-|A_0|}\) gives exactly the displayed formula in the
statement.  All real factors are nonnegative because \(0\le p\le1\) and
\(0\le t_v\le1\), so the uses of \(\operatorname{ofReal}\) and finite
additivity above take place inside finite real masses and then translate back
to extended nonnegative real measure values.
\end{proof}
